\documentclass[11pt]{article}

\usepackage[margin=1in]{geometry}
\usepackage{amsmath}
\usepackage{booktabs}
\usepackage[hidelinks]{hyperref}
\usepackage{microtype}

\title{Bulk-Silicon Parent Qualification and\\Joint-Reduction Technical Report}
\author{Project Koios}
\date{Draft calculation report}

\begin{document}
\maketitle

\begin{abstract}
This report defines the calculation boundary for a prospective
scalar-relativistic, non-spin-polarized, non-spin--orbit-coupled bulk-silicon
parent, a compatible ten-spatial-orbital Wannier bridge, and a finite-support
local-model hierarchy.  Existing evidence is limited to an audited provisional
direct bootstrap.  No production parent, qualified Wannier representation,
joint-fidelity result, or material-validation claim is established.
\end{abstract}

\section{Evidence status}

This is a publication-independent calculation report.  It describes a prospective
material calculation with an audited provisional bootstrap.  It is not a completed
production calculation, validated silicon Hamiltonian, qualified Wannier
representation, model-selection result, or scientific validation.  No new scientific
calculator was executed for this report.

\section{Scientific objective}

The target is a pristine bulk-silicon parent represented consistently at three levels:
\begin{enumerate}
  \item a first-principles plane-wave oracle $O_A(\mathbf{k})$;
  \item a retained ten-spatial-orbital Wannier bridge $O_B(\mathbf{k})$; and
  \item a finite-support $sp^3s^*$ local-model hierarchy
    $O_C^{U,S}(\mathbf{k})$.
\end{enumerate}

For retained frame $Q(\mathbf{k})$, alignment $U(\mathbf{k})$, support $S$, and
real-space blocks $T_U(\mathbf{R})$, the intended relations are
\begin{align}
  O_B(\mathbf{k})
    &= Q(\mathbf{k})^\dagger O_A(\mathbf{k})Q(\mathbf{k}), \\
  O_C^{U,S}(\mathbf{k})
    &= \sum_{\mathbf{R}\in S}
       e^{i\mathbf{k}\cdot\mathbf{R}}T_U(\mathbf{R}).
\end{align}
A candidate is adequate only if the same model, support, alignment, basis, and energy
reference satisfy both declared spectral and aligned-operator tolerances.  Visual band
agreement alone is insufficient.

\section{External empirical development control}

The calculation-local Vogl $sp^3s^*$ control follows the empirical model of
Vogl, Hjalmarson, and Dow~\cite{Vogl1983} and contains the accepted complete Table~1
transcription: 16 compounds, 13 parameters per compound, and energies in eV.  An
independent implementation constructs the nearest-neighbor, spinless, non-SOC ten-band
Bloch Hamiltonian in the ordered anion--cation basis
$(s,p_x,p_y,p_z,s^*)_a\oplus(s,p_x,p_y,p_z,s^*)_c$.  Its reduced Cartesian wavevector
uses $\mathbf{k}=(2\pi/a_L)\mathbf{q}$.

Software and analytic checks cover the table inventory, silicon parameter row,
structure factors at symmetry points, exact Hermitian assembly, the silicon $\Gamma$
spectrum, reciprocal periodicity, and invalid-input rejection.  This external empirical
control is not derived from the provisional first-principles bootstrap or a Wannier
bridge and does not freeze the candidate inventory, establish model adequacy, or qualify
a material result.

\section{Existing evidence}

The historical direct bootstrap evaluated six wavefunction-cutoff candidates and four
reciprocal-mesh candidates, with exact reuse of the C48 branch as logical K8.  This
produced nine unique self-consistent-field (SCF) and nine linked diagnostic
non-self-consistent-field (NSCF) executions.  The compact audit recomputed the required
identities and observables and found the retained matrix complete and internally
consistent.

Several adjacent comparisons meet the frozen finite-setting criteria.  Human
disposition selected C48/K8---48/192~Ry with a shifted $8\times8\times8$ mesh---as a
provisional finite-setting parameter pair.  This establishes neither an accepted
production setting nor an infinite-basis error bound.  The provisional geometry is not
an accepted equilibrium lattice, the three diagnostic momenta do not locate the
conduction valley, and printed fixed-point bands do not converge an effective mass.

Every invocation also emitted the same IEEE flag report.  Instrumented diagnoses
localized one event to external OpenCL discovery during startup and classified a
separate cross-BLAS, cross-MPI divide-by-zero below \texttt{ZHEGVX} during
initial-wavefunction subspace diagonalization.  The effect on retained observables and
the sources of the remaining reported flag classes are unresolved.  Successful exit
and \texttt{JOB DONE.} do not establish harmlessness.  A mixed cutoff--mesh interaction
check implied by C48/K8 would be new protected execution.

Exact values, hashes, criteria, and warning investigations remain at the historical
locations catalogued in
\path{calculations/Si/bulk/evidence-crosswalk.md}.  They are referenced rather than
copied or relabelled.

\section{Required parent qualification}

The production parent must freeze the physical branch, lattice, PBE/ONCV approximation,
Quantum ESPRESSO executable, pseudopotential identity, occupations, solver settings,
cutoff and mesh inventory, compact/native artifact boundary, and warning policy before
execution.  Qualification must keep energy, stress, fixed-point bands, gaps, valley
identity, and curvature requirements distinct.

For each retained observable $q$, finite-setting agreement
\begin{equation}
  \lvert q_j-q_{j+1}\rvert \leq \tau_q
\end{equation}
demonstrates only stability across the tested pair.  It does not imply a bound on
$\lvert q_j-q_\infty\rvert$ without an additional reviewed error argument.

\section{Required retained-space bridge}

The Wannier bridge must retain projections, windows, centers, spreads, convergence
status, reciprocal sewing, symmetry and orbital semantics, energy reference, and direct
parent comparisons.  At crossings or degeneracies, comparison must use spectral
clusters or projectors rather than undocumented individual eigenvector labels.

Wannier90 may produce a candidate witness.  Wannier90 convergence is not by itself a
proof of the bridge contract.  The tutorial-scale silicon Wannier extraction remains an
integration precursor and is excluded from the production admissible sets.

\section{Joint spectral/operator decision}

For a frozen finite candidate inventory $D$, define
\begin{equation}
  F(\tau_E,\tau_O)
  = \left\{d\in D:
      L_E(d)\leq\tau_E\ \land\ L_O(d)\leq\tau_O
    \right\}.
\end{equation}
If this set is nonempty, the intended finite decision is
\begin{equation}
  C^*(\tau_E,\tau_O)
  = \min_{d\in F(\tau_E,\tau_O)} C(d).
\end{equation}
The same $d$ must determine the model, support, alignment, and energy reference in both
loss channels.  Separate spectral and operator optima do not form a common witness.

The decision must preserve distinct numerical, subspace, alignment, interpolation, and
support-truncation error channels.  It must report one of the following:
\begin{itemize}
  \item a common witness and the smallest adequate declared model;
  \item a reviewed separation result for an exhaustively evaluated finite inventory;
    or
  \item an explicit inconclusive outcome.
\end{itemize}
Failure of a local optimizer does not prove infeasibility.

\section{Withheld analysis and residual attribution}

Construction and withheld momentum sets must be disjoint by contract.  Withheld
analysis must report band, gap, valley, effective-mass, projector, represented-operator,
and locality channels separately.  Stable residuals may be attributed by onsite or
hopping block, orbital semantics, symmetry channel, and support shell only when the
representation and alignment contracts permit that interpretation.

Sensitivity analysis must include every prospectively declared parent, Wannier,
stencil, alignment-map resolution, and threshold variation.  Post-hoc sensitivity does
not retroactively justify thresholds.

\section{Verification and retention}

A completed calculation requires independent verification of sources, schemas,
identities, finite candidate coverage, objective reconstruction, threshold evaluation,
and withheld results.  Generated plots must consume verified records.  A retained
package must be no-overwrite and checksum-bound.

Formal proofs or certificates may establish exact implications for finite definitions
and accepted interval premises.  They do not certify Quantum ESPRESSO or Wannier90
execution, JSON, binary64 rounding, pseudopotential transferability, or physical
validation.

\section{Current gate and next decision}

The active successor disposition permits explicitly labelled downstream development
and exploratory analysis using provisional C48/K8 evidence.  Production-parent and
publication-grade acceptance remain pending an observable-level comparison through an
independently identified established QE/toolchain path and the absent mixed
cutoff--mesh interaction check.  Capturing and replaying the exact first
$8\times8$ subspace Hamiltonian and overlap matrices remains an optional focused
diagnosis rather than a prerequisite for all downstream work.  Every new QE invocation
requires separate authorization.  The detailed gates are recorded in
\path{calculations/Si/bulk/development-plan.md} and the active decision in
\path{calculations/Si/bulk/provisional-use-disposition.md}.

\section{Nonclaims}

This report does not establish production convergence, an infinite-basis result,
experimental agreement, pseudopotential transferability, a qualified Wannier model,
an adequate tight-binding model, or scientific validation.  It authorizes no Quantum
ESPRESSO or Wannier90 execution, evidence promotion, publication, or release.

\appendix
\section{Supplemental material}

\subsection{Technical note: Vogl-control implementation provenance}

The scientific definition and parameter authority for the empirical control are the
published Vogl model and its Table~1~\cite{Vogl1983}.  A separate development lineage
was inspected for comparison: a public zinc-blende LCAO Gist by
Mahdisadjadi~\cite{MahdisadjadiGist} and the 2026
\texttt{SiliconLCAO\_4sDOS-v3} notebook attributed within the notebook to Simon
Phillpot~\cite{Phillpot2026}.  The notebook states that it began from the Gist, corrected
phase factors, and added the $s^*$ orbitals, with modifications dated 29~April through
11~May 2026.

The calculation-local module
\path{vogl1983_lcao_sp3sstar/model.py} was written independently from the published
model equations and the accepted Table~1 transcription.  No Gist or notebook source code
was copied.  Their role is limited to development provenance and numerical cross-checking;
they do not supersede the primary source.  The notebook agreed with 203 of the 208 Table~1
cells at four-decimal numerical normalization.  The five disagreements were resolved in
favor of the printed Vogl values.  Neither the Gist nor the notebook exposed an explicit
license in the inspected material, which is an additional reason to keep their role
separate from the maintained implementation.

\subsection{Technical note: LCAO and tetrahedral \texorpdfstring{$sp^3$}{sp3} construction}

This summary follows the tight-binding treatments of Harrison, Chadi and Cohen, and
Vogl et al.~\cite{Harrison1980,ChadiCohen1975,Vogl1983}.  At one site, begin with the
orthonormal orbitals $\lvert s\rangle$, $\lvert p_x\rangle$, $\lvert p_y\rangle$, and
$\lvert p_z\rangle$,
\begin{equation}
  \langle \alpha\mid\beta\rangle=\delta_{\alpha\beta},
  \qquad \alpha,\beta\in\{s,p_x,p_y,p_z\}.
\end{equation}
The four $sp^3$ hybrid orbitals are formed by combining the $s$, $p_x$, $p_y$,
and $p_z$ orbitals with equal magnitudes but different sign patterns:
\begin{equation}
  \begin{pmatrix}
    \lvert h_1\rangle \\
    \lvert h_2\rangle \\
    \lvert h_3\rangle \\
    \lvert h_4\rangle
  \end{pmatrix}
  = \frac12
  \underbrace{
  \begin{pmatrix}
     1 &  1 &  1 &  1 \\
     1 &  1 & -1 & -1 \\
     1 & -1 &  1 & -1 \\
     1 & -1 & -1 &  1
  \end{pmatrix}}_{H}
  \begin{pmatrix}
    \lvert s\rangle \\
    \lvert p_x\rangle \\
    \lvert p_y\rangle \\
    \lvert p_z\rangle
  \end{pmatrix}.
\end{equation}
The sign matrix $H$ is a $4\times4$ Hadamard matrix.  Its rows are mutually
orthogonal and satisfy
\begin{equation}
  HH^{\mathsf{T}}=4I.
\end{equation}
Consequently, $H/2$ defines an orthogonal change of basis that preserves norms and
inner products.  It transforms the orthonormal atomic-orbital basis into four
orthonormal hybrid orbitals, $\langle h_m\mid h_n\rangle=\delta_{mn}$.  Each row's sign
pattern selects one of the $(1,1,1)$, $(1,\bar{1},\bar{1})$,
$(\bar{1},1,\bar{1})$, or $(\bar{1},\bar{1},1)$ tetrahedral bond directions.  Because
this normalized Hadamard transform is symmetric and orthogonal, the same matrix also
recovers the original $s$ and $p$ orbitals from the hybrid basis.

For a periodic solid, form Bloch sums of localized orbitals,
\begin{equation}
  \lvert \alpha b\mathbf{k}\rangle
  = \frac{1}{\sqrt{N}}\sum_{\mathbf{R}}
    e^{i\mathbf{k}\cdot(\mathbf{R}+\boldsymbol{\nu}_b)}
    \lvert \alpha b\mathbf{R}\rangle,
\end{equation}
where $b$ labels the anion or cation sublattice.  Translation therefore enters a
Hamiltonian matrix element through a phase $e^{i\mathbf{k}\cdot\mathbf{r}}$; reversing
the bond gives the complex-conjugate phase required by Hermiticity.  The LCAO
approximation represents each symmetry-distinct localized matrix element by a constant:
onsite $E_s$ and $E_p$, and nearest-neighbor $V_{ss}$, $V_{sp}$, $V_{xx}$, and
$V_{xy}$.  In a heteropolar crystal, anion and cation onsite energies differ, and
$V(s_a,p_c)$ and $V(s_c,p_a)$ are separate directional parameters.  Onsite $s$--$p$
matrix elements vanish by the site symmetry of the orthogonal basis; intersite $s$--$p$
entries acquire signs from the odd parity of the $p$ orbital.

For the four nearest-neighbor displacements
\begin{equation}
  \mathbf{r}_{1,2,3,4}
  = \frac{a_L}{4}
    \bigl\{(1,1,1),(1,-1,-1),(-1,1,-1),(-1,-1,1)\bigr\},
\end{equation}
write $e_n=e^{i\mathbf{k}\cdot\mathbf{r}_n}$.  The four structure factors are
\begin{align}
  g_0 &= \tfrac14(e_1+e_2+e_3+e_4), &
  g_1 &= \tfrac14(e_1+e_2-e_3-e_4), \\
  g_2 &= \tfrac14(e_1-e_2+e_3-e_4), &
  g_3 &= \tfrac14(e_1-e_2-e_3+e_4).
\end{align}
Thus representative Bloch-matrix entries are $g_0V(s,s)$ for $s$--$s$ coupling and
$g_jV(s_a,p_{j,c})$ for $s$--$p_j$ coupling, with the reverse entries fixed by the
Hermitian conjugate and the orbital-direction sign convention.  The $sp^3$ construction
produces an $8\times8$ Hamiltonian for two four-orbital sites.  Vogl's $sp^3s^*$ model
adds one excited $s^*$ orbital per site and only the declared nearest-neighbor
$s^*$--$p$ couplings, producing the maintained $10\times10$ block form
\begin{equation}
  H(\mathbf{k})=
  \begin{pmatrix}
    D_a & T(\mathbf{k}) \\
    T(\mathbf{k})^\dagger & D_c
  \end{pmatrix},
  \qquad
  D_b=\operatorname{diag}
  \bigl(E(s,b),E(p,b),E(p,b),E(p,b),E(s^*,b)\bigr).
\end{equation}
The Table~1 transfer parameters determine $T(\mathbf{k})$; the $s^*$--$s$ and
$s^*$--$s^*$ transfers are structural zeros in this model.

\bibliographystyle{unsrt}
\bibliography{references}

\end{document}
